\documentclass[12pt,a4paper]{report}

% --- Packages ---
\usepackage[utf8]{inputenc}
\usepackage[T1]{fontenc}
\usepackage[vietnamese,english]{babel}
\usepackage{graphicx}
\usepackage{amsmath}
\usepackage{hyperref}
\usepackage{float}
\usepackage{geometry}
\usepackage{booktabs}
\usepackage{enumitem}
\usepackage{caption}

% --- Setup ---
\geometry{margin=1in}
\hypersetup{
    colorlinks=true,
    linkcolor=blue,
    filecolor=magenta,      
    urlcolor=cyan,
    pdftitle={HanoiGO - Module 1 Documentation},
}

\begin{document}

\title{HanoiGO Technical Documentation \\ \large Module 1: Trip Planner}
\author{Nguyen Hung}
\date{\today}
\maketitle

\chapter{Module 1: Trip Planner -- Intelligent Itinerary Generation}
\label{chap:trip-planner}

\section{Overview}

The Trip Planner module is the core scheduling engine of HanoiGO. It receives a list of tourist attractions selected by the user, along with constraints (number of days, daily time window, travel date, user GPS location), and produces an optimized multi-day itinerary that minimizes total travel time and waiting time while respecting each place's operating hours.

\subsection{Technology Stack}

\begin{table}[H]
\centering
\begin{tabular}{|l|l|}
\hline
\textbf{Layer} & \textbf{Technology} \\
\hline
Backend Framework & NestJS (TypeScript) \\
Database ORM & Prisma (PostgreSQL) \\
Distance Matrix API & Goong Maps API (with Haversine fallback) \\
Frontend Framework & Next.js 14 (React, TypeScript) \\
Map Rendering & Leaflet + CartoDB Voyager tiles \\
Routing Visualization & Goong Directions API \\
\hline
\end{tabular}
\caption{Technology stack for the Trip Planner module}
\end{table}

\subsection{System Architecture}

\begin{figure}[H]
\centering
\begin{verbatim}
+-------------------+       POST /trips/generate-itinerary       +------------------+
|   Next.js Client  | -------------------------------------->    |  NestJS Backend  |
|  (DiscoveryPage)  | <--------------------------------------   | (TripPlanner     |
|                   |           JSON ItineraryResponse           |    Service)      |
+-------------------+                                            +------------------+
        |                                                               |
        | Leaflet GPS                                                   | Prisma
        | (onLocationFound)                                             | $queryRawUnsafe
        v                                                               v
  [Browser GPS]                                                   [PostgreSQL]
                                                                        |
                                                              Goong DistanceMatrix API
                                                              (or Haversine fallback)
\end{verbatim}
\caption{High-level architecture of Module 1}
\end{figure}

%% =====================================================================
\section{Algorithm Pipeline}
\label{sec:algorithm-pipeline}

The itinerary generation follows a 6-step pipeline executed server-side within \texttt{TripPlannerService.generateItinerary()}:

\begin{enumerate}
    \item \textbf{Data Fetching} -- Query the \texttt{places} table by name.
    \item \textbf{Pre-Filter} -- Remove places closed on all travel days.
    \item \textbf{K-Means Clustering} -- Group places geographically into $k$ day-clusters.
    \item \textbf{Greedy Nearest Neighbor (GNN) with Time Windows} -- Sequence each cluster optimally.
    \item \textbf{Conflict Resolution} -- Reassign dropped places to days with remaining capacity.
    \item \textbf{Response Formatting} -- Convert internal structures to the client-facing JSON.
\end{enumerate}

%% ── Step 0 ──────────────────────────────────────────────────────────
\subsection{Step 0: Data Fetching}

The service receives a \texttt{GenerateItineraryDto} containing:

\begin{table}[H]
\centering
\begin{tabular}{|l|l|l|}
\hline
\textbf{Field} & \textbf{Type} & \textbf{Example} \\
\hline
\texttt{placeNames} & \texttt{string[]} & \texttt{["Hoa Lo Prison", "St. Joseph's Cathedral"]} \\
\texttt{numDays} & \texttt{number} & \texttt{2} \\
\texttt{startTime} & \texttt{number} & \texttt{480} (08:00) \\
\texttt{endTime} & \texttt{number} & \texttt{1080} (18:00) \\
\texttt{travelDate} & \texttt{string} & \texttt{"2026-05-01"} \\
\texttt{visitDurationMin} & \texttt{number} & \texttt{30} \\
\texttt{startLat} & \texttt{number?} & \texttt{21.0242} \\
\texttt{startLng} & \texttt{number?} & \texttt{105.8576} \\
\hline
\end{tabular}
\caption{Input DTO fields}
\end{table}

A raw SQL query fetches matching places from PostgreSQL:
\begin{verbatim}
SELECT id, name, category, district, lat, lng, image_url,
       always_open, open_days, open_time_start, open_time_end,
       has_break, break_start, break_end, visit_duration_min
FROM places WHERE name = ANY($1)
\end{verbatim}

Each database record is mapped to an internal \texttt{Place} interface. Time fields (\texttt{open\_time\_start}, etc.) are converted from \texttt{TIME} type to minutes-from-midnight using the \texttt{dbTimeToMin()} utility.

%% ── Step 1 ──────────────────────────────────────────────────────────
\subsection{Step 1: Pre-Filter}

For each place, the algorithm checks whether \texttt{openDays} intersects with any of the travel days $[d_0, d_0+1, \ldots, d_0+k-1]$. Places that are \texttt{alwaysOpen = true} pass automatically. Places closed on \emph{all} travel days are moved to the \texttt{infeasible[]} output array with a human-readable reason string.

\textbf{Complexity:} $O(n \cdot k)$ where $n$ = number of places, $k$ = number of days.

%% ── Step 2 ──────────────────────────────────────────────────────────
\subsection{Step 2: K-Means Clustering}

The algorithm partitions feasible places into $k$ geographic clusters (one per day) using Lloyd's K-Means algorithm with the Haversine distance metric.

\subsubsection{Haversine Distance}
\begin{equation}
d = 2R \cdot \arcsin\left(\sqrt{\sin^2\!\left(\frac{\Delta\phi}{2}\right) + \cos\phi_1 \cdot \cos\phi_2 \cdot \sin^2\!\left(\frac{\Delta\lambda}{2}\right)}\right)
\end{equation}
where $R = 6371$ km (Earth radius), $\phi$ = latitude in radians, $\lambda$ = longitude in radians.

\subsubsection{Centroid Initialization}
Centroids are initialized by uniform sampling: $c_i = \text{places}[i \cdot \lfloor n/k \rfloor]$ for $i \in [0, k)$.

\subsubsection{Convergence}
The algorithm iterates until centroids shift by less than $0.0001\degree$ in both latitude and longitude, or a maximum of 50 iterations is reached.

\subsubsection{Rebalancing}
After convergence, clusters exceeding \texttt{MAX\_PLACES\_PER\_DAY = 5} are rebalanced. The place closest to the lightest cluster's centroid (by the score $d_{\text{dest}} - d_{\text{source}}$) is moved iteratively until all clusters satisfy the capacity constraint.

%% ── Step 3 ──────────────────────────────────────────────────────────
\subsection{Step 3: Distance Matrix}

For each cluster, a pairwise travel-duration matrix $M[i][j]$ (in seconds) is computed.

\begin{itemize}
    \item \textbf{Primary:} Goong Maps DistanceMatrix API (\texttt{vehicle=bike}).
    \item \textbf{Fallback:} If the API key is missing, rate-limited, or returns an error, the system falls back to Haversine-based estimation at an average speed of 15~km/h.
    \item \textbf{Rate-limit handling:} On \texttt{OVER\_RATE\_LIMIT}, the system waits 3 seconds and retries once before falling back.
\end{itemize}

%% ── Step 4 ──────────────────────────────────────────────────────────
\subsection{Step 4: Greedy Nearest Neighbor with Time Windows}

This is the core scheduling algorithm for each day. It selects the visit order within a cluster.

\subsubsection{First Stop Selection (Start Location Optimization)}
If the user provides \texttt{startLat}/\texttt{startLng}, the algorithm finds the place closest to the user's current GPS location using Haversine distance and uses it as the first stop. Otherwise, it defaults to index 0.

\subsubsection{Greedy Loop}
From the current stop, the algorithm evaluates all unvisited places and selects the one with the minimum \emph{cost}:
\begin{equation}
\text{cost}(j) = \text{travelTime}_{i \to j} + \text{waitTime}(j) \times 60
\end{equation}

A candidate $j$ is \textbf{skipped} if:
\begin{itemize}
    \item The place is closed on the current day of week.
    \item Arrival time exceeds \texttt{effectiveClose $-$ visitDuration} (would arrive too late).
    \item Departure time would exceed the user's \texttt{endTime}.
\end{itemize}

\textbf{Wait time} is calculated for two scenarios:
\begin{itemize}
    \item Arriving before \texttt{openTimeStart}: $\text{wait} = \text{openTimeStart} - \text{arriveMin}$
    \item Arriving during a break period: $\text{wait} = \text{breakEnd} - \text{arriveMin}$
\end{itemize}

\subsubsection{Time Cascade}
After the GNN produces a route, if \texttt{startLat}/\texttt{startLng} is provided, the travel time from the user's location to the first stop is computed and all subsequent stop times are cascaded (shifted forward) accordingly.

\textbf{Complexity:} $O(m^2)$ per cluster where $m$ = cluster size (at most 5).

%% ── Step 5 ──────────────────────────────────────────────────────────
\subsection{Step 5: Conflict Resolution}

Places that were dropped by GNN (could not fit within the time window) are attempted to be reassigned to other days that have sufficient remaining time:

\begin{equation}
\text{availableTime} = \text{endTime} - \text{lastStop.departMin}
\end{equation}

A place is reassigned if $\text{availableTime} \geq \text{visitDuration} + \text{PARKING\_BUFFER} + 15$. Places that cannot fit in any day are reported in the \texttt{unscheduled[]} output array.

%% ── Step 6 ──────────────────────────────────────────────────────────
\subsection{Step 6: Response Formatting}

The internal \texttt{DayItinerary[]} is transformed into the client-facing \texttt{ItineraryResponse} JSON:

\begin{verbatim}
{
  "days": [{
    "dayNumber": 1,
    "dayLabel": "Fri",
    "color": "#3B82F6",
    "stops": [{
      "order": 1,
      "name": "St. Joseph's Cathedral",
      "arriveAt": "08:15",
      "departAt": "08:45",
      "travelFromPrevMin": 5,
      "waitMin": 0
    }, ...],
    "totalTravelMin": 32,
    "freeTimeMin": 408
  }],
  "infeasible": [],
  "unscheduled": []
}
\end{verbatim}

%% =====================================================================
\section{Frontend Integration}

\subsection{User Flow}

\begin{enumerate}
    \item User navigates to \texttt{/discovery} page.
    \item The Leaflet map initializes and begins GPS tracking (blue dot).
    \item User clicks \textbf{``Lên kế hoạch chuyến đi''} to enter Plan Mode.
    \item User selects 2+ places from the card list.
    \item User clicks \textbf{``Tối ưu lịch trình''} to open the configuration modal.
    \item User configures: number of days, travel date, start/end time, visit duration.
    \item User clicks \textbf{``Tạo lịch trình''}.
    \item Frontend sends \texttt{POST /trips/generate-itinerary} with the cached GPS coordinates.
    \item The sidebar switches to the Timeline view showing the optimized itinerary.
    \item The map displays numbered markers with day-specific colors and polyline routes.
\end{enumerate}

\subsection{GPS Coordinate Flow}

The GPS location is obtained from Leaflet's \texttt{map.locate()} and propagated upward via the \texttt{onLocationFound} callback from \texttt{DiscoveryMap} to \texttt{DiscoveryPage}. This cached coordinate is sent as \texttt{startLat}/\texttt{startLng} in the API request, ensuring the backend always receives the user's position without relying on a separate \texttt{getCurrentPosition} call at generation time.

%% =====================================================================
\section{Logic Review Summary}

After a thorough code review, the following aspects have been verified:

\begin{table}[H]
\centering
\begin{tabular}{|l|l|l|}
\hline
\textbf{Aspect} & \textbf{Status} & \textbf{Notes} \\
\hline
Pre-filter logic & \checkmark OK & Correctly checks openDays against all travel days \\
K-Means convergence & \checkmark OK & Convergence threshold $0.0001\degree$, max 50 iters \\
K-Means rebalancing & \checkmark OK & Moves places to lightest cluster by score \\
Empty cluster handling & \checkmark OK & Filtered out after convergence \\
Goong API fallback & \checkmark OK & Haversine fallback on error/missing key \\
Rate-limit retry & \checkmark OK & Single retry after 3s sleep \\
GNN first-stop selection & \checkmark OK & Uses Haversine from startLat/startLng \\
GNN time-window cost & \checkmark OK & cost = travelSec + waitMin * 60 \\
Break-time handling & \checkmark OK & Detects arrival during break, adds wait \\
Time cascade after GPS & \checkmark OK & Shifts all stops after first-stop recalc \\
Conflict resolution & \checkmark OK & Tries all days, checks openDays compat \\
Frontend GPS propagation & \checkmark OK & Cached from Leaflet, not re-requested \\
\hline
\end{tabular}
\caption{Logic review checklist}
\end{table}

\subsection{Known Limitations}

\begin{enumerate}
    \item \textbf{GNN is greedy, not globally optimal.} It does not guarantee the TSP-optimal route. For $\leq 5$ places per day, the difference is negligible.
    \item \textbf{K-Means initialization is deterministic} (uniform sampling), which may produce suboptimal clusters for certain spatial distributions. K-Means++ could improve this.
    \item \textbf{Conflict resolution appends to the end of a day} using a fixed 10-minute travel estimate (\texttt{travelFromPrevSec = 600}), which is an approximation.
    \item \textbf{No return-to-origin cost.} The algorithm does not consider travel time back to the user's starting point at day's end.
\end{enumerate}

%% =====================================================================
\section{Test Cases}
\label{sec:test-cases}

All test cases use Jest with NestJS \texttt{TestingModule}. The Prisma database is mocked with \texttt{jest.fn()} to return predefined place data. All 6 tests pass successfully.

\subsection{Test Data Setup}

\begin{table}[H]
\centering
\small
\begin{tabular}{|l|l|l|l|l|}
\hline
\textbf{ID} & \textbf{Name} & \textbf{District} & \textbf{Lat, Lng} & \textbf{Special} \\
\hline
hk1 & Hồ Hoàn Kiếm & Hoàn Kiếm & 21.0285, 105.8542 & alwaysOpen \\
hk2 & Nhà Thờ Lớn & Hoàn Kiếm & 21.0287, 105.8489 & alwaysOpen \\
hk3 & Phố Cổ Hà Nội & Hoàn Kiếm & 21.0345, 105.8516 & alwaysOpen \\
cg1 & Bảo tàng Dân tộc học & Cầu Giấy & 21.0404, 105.7955 & Closed Mon \\
cg2 & Công viên Cầu Giấy & Cầu Giấy & 21.0264, 105.7938 & alwaysOpen \\
cg3 & Quán Cafe Cầu Giấy & Cầu Giấy & 21.0280, 105.7950 & Opens 14:00 \\
\hline
\end{tabular}
\caption{Mock place data for unit tests}
\end{table}

%% ── TC1 ─────────────────────────────────────────────────────────────
\subsection{TC-01: K-Means Spatial Clustering}

\begin{table}[H]
\centering
\begin{tabular}{|p{3cm}|p{10cm}|}
\hline
\textbf{Objective} & Verify that geographically distant places are separated into different day-clusters. \\
\hline
\textbf{Input} & 5 places: 3 in Hoàn Kiếm + 2 in Cầu Giấy. \texttt{numDays = 2}. \\
\hline
\textbf{Expected} & Day 1 contains only Hoàn Kiếm places. Day 2 contains only Cầu Giấy places (or vice versa). No cross-district mixing. \\
\hline
\textbf{Result} & \checkmark PASS \\
\hline
\end{tabular}
\end{table}

%% ── TC2 ─────────────────────────────────────────────────────────────
\subsection{TC-02: Pre-Filter -- Closed Day Exclusion}

\begin{table}[H]
\centering
\begin{tabular}{|p{3cm}|p{10cm}|}
\hline
\textbf{Objective} & Verify that places closed on the travel date are excluded from scheduling and reported as infeasible. \\
\hline
\textbf{Input} & 2 places: Hồ Hoàn Kiếm (alwaysOpen) + Bảo tàng DTH (closed Monday). \texttt{travelDate = Monday}. \texttt{numDays = 1}. \\
\hline
\textbf{Expected} & Bảo tàng DTH appears in \texttt{infeasible[]} array. Hồ Hoàn Kiếm is scheduled normally. \\
\hline
\textbf{Result} & \checkmark PASS \\
\hline
\end{tabular}
\end{table}

%% ── TC3 ─────────────────────────────────────────────────────────────
\subsection{TC-03: Time Window Optimization}

\begin{table}[H]
\centering
\begin{tabular}{|p{3cm}|p{10cm}|}
\hline
\textbf{Objective} & Verify that the GNN algorithm respects opening hours and correctly calculates wait time for places that open later in the day. \\
\hline
\textbf{Input} & 2 places in same area: Công viên (alwaysOpen) + Quán Cafe (opens 14:00). \texttt{startTime = 08:00}. \\
\hline
\textbf{Expected} & Công viên is scheduled first (morning). Quán Cafe is scheduled second. \texttt{waitMin > 0} for the cafe stop since arrival precedes its 14:00 opening. \\
\hline
\textbf{Result} & \checkmark PASS \\
\hline
\end{tabular}
\end{table}

%% ── TC4 ─────────────────────────────────────────────────────────────
\subsection{TC-04: Start Location Optimization (Case A)}

\begin{table}[H]
\centering
\begin{tabular}{|p{3cm}|p{10cm}|}
\hline
\textbf{Objective} & Verify that when the user provides GPS coordinates, the algorithm selects the nearest place as the first stop. \\
\hline
\textbf{Input} & 2 places: Hồ Hoàn Kiếm + Phố Cổ. User at Nhà Hát Lớn (21.0242, 105.8576) -- closer to Hồ Hoàn Kiếm. \\
\hline
\textbf{Expected} & Stop \#1 = Hồ Hoàn Kiếm. Stop \#2 = Phố Cổ Hà Nội. \\
\hline
\textbf{Result} & \checkmark PASS \\
\hline
\end{tabular}
\end{table}

%% ── TC5 ─────────────────────────────────────────────────────────────
\subsection{TC-05: Start Location Optimization (Case B -- Reversed)}

\begin{table}[H]
\centering
\begin{tabular}{|p{3cm}|p{10cm}|}
\hline
\textbf{Objective} & Confirm that changing the user's start location reverses the stop order correctly. \\
\hline
\textbf{Input} & Same 2 places. User at Chợ Đồng Xuân (21.0378, 105.8499) -- closer to Phố Cổ. \\
\hline
\textbf{Expected} & Stop \#1 = Phố Cổ Hà Nội. Stop \#2 = Hồ Hoàn Kiếm. \\
\hline
\textbf{Result} & \checkmark PASS \\
\hline
\end{tabular}
\end{table}

%% ── TC6 ─────────────────────────────────────────────────────────────
\subsection{TC-06: Overloaded Time Budget}

\begin{table}[H]
\centering
\begin{tabular}{|p{3cm}|p{10cm}|}
\hline
\textbf{Objective} & Verify that excess places are gracefully dropped when the daily time budget is insufficient. \\
\hline
\textbf{Input} & 6 places, each 60 min visit. \texttt{numDays = 1}, \texttt{startTime = 08:00}, \texttt{endTime = 13:00} (5-hour budget). Total visit time alone = 6 hours. \\
\hline
\textbf{Expected} & 3--4 places scheduled in \texttt{days[0].stops}. Remaining places in \texttt{unscheduled[]} with reason. \\
\hline
\textbf{Result} & Expected behavior (manually verified). \\
\hline
\end{tabular}
\end{table}

%% ── TC7 ─────────────────────────────────────────────────────────────
\subsection{TC-07: Goong API Fallback}

\begin{table}[H]
\centering
\begin{tabular}{|p{3cm}|p{10cm}|}
\hline
\textbf{Objective} & Verify the system does not crash when the Goong API is unavailable. \\
\hline
\textbf{Input} & Set \texttt{GOONG\_API\_KEY} to an invalid value. Request itinerary for 3+ places. \\
\hline
\textbf{Expected} & System logs a warning, falls back to Haversine-based distance estimation (15 km/h average), and still returns a valid itinerary. \\
\hline
\textbf{Result} & Expected behavior (manually verified via log inspection). \\
\hline
\end{tabular}
\end{table}

\section{Conclusion}

The Trip Planner module implements a practical heuristic pipeline combining K-Means spatial clustering with a Greedy Nearest Neighbor scheduling algorithm enhanced by time-window awareness. The start-location optimization ensures the first stop is always the place nearest to the user's current GPS position, providing a natural and efficient travel experience. The system gracefully handles edge cases including closed days, break periods, API failures, and time budget overflows.

\end{document}
